\chapter{Hysterectomy}

\section*{Overview}
A hysterectomy is surgical removal of the uterus. Depending on the condition, the cervix, fallopian tubes, and ovaries may also be removed.

\section*{Why This Test or Procedure Is Ordered}
It may be performed for symptomatic fibroids, abnormal uterine bleeding, endometriosis, uterine prolapse, cancer, or conditions that have not responded to less invasive treatment.

\section*{How This Test or Procedure Is Performed}
The uterus may be removed through an abdominal incision, laparoscopically or robotically, or through the vagina. The extent of surgery depends on the diagnosis and reproductive or oncologic needs.

\section*{Risks and Recovery}
Bleeding, infection, blood clots, injury to the bladder, bowel, or ureters, anesthesia complications, and early menopause when ovaries are removed are possible. Recovery varies by approach; pregnancy is not possible after hysterectomy.

\section*{References}
\begin{enumerate}
  \item MedlinePlus. Hysterectomy. U.S. National Library of Medicine; 2025.
  \item Current Medical Diagnosis and Treatment. McGraw-Hill; 2025.
\end{enumerate}
\clearpage
